
\documentclass[10pt]{article}
\usepackage{amsmath, amssymb, bm}
\usepackage[a4paper, margin=0.8in]{geometry}
\begin{document}
\section*{GFEST: A compact mathematical sketch}
We posit a response $y\in\mathbb{R}$ and inputs $\mathbf{x}\in\mathbb{R}^d$.
GFEST searches a restricted, interpretable hypothesis class
\begin{align}
\mathcal{H}=\left\{\, h(\mathbf{x})=\beta_0+\sum_{t=1}^K \beta_t\, \phi_t(\mathbf{x}) \ \middle|\ \phi_t\in\Phi \right\},
\end{align}
where $\Phi$ is built from protected primitives (e.g., $\mathrm{id}, \mathrm{square}, \mathrm{sqrt}, \log(1{+}|x|), \exp, \sin, \cos, 1/x$) applied to one or two variables (products allowed).

\paragraph{Chromosome (fixed-length).} Each $\phi_t$ is encoded by a tuple $(\text{family}, j, k, p, q, a)$ where \emph{family}$\in\{\text{unary}, \text{product}\}$, $j,k$ are variable indices, $p,q$ select primitives, and $a\in\{0,1\}$ toggles activity.

\paragraph{Fitting.} For a fixed set of active $\{\phi_t\}$, we fit $\bm{\beta}$ via OLS or LASSO:
\begin{align}
\widehat{\bm{\beta}}=\arg\min_{\bm{\beta}}\ \frac{1}{n}\sum_{i=1}^n\big(y_i-\beta_0-\sum_t\beta_t\phi_t(\mathbf{x}_i)\big)^2 + \lambda\|\bm{\beta}\|_1.
\end{align}

\paragraph{Model selection.} Assuming homoscedastic Gaussian noise, the AICc objective reads
\begin{align}
\mathrm{AICc}=n\ln(\widehat{\sigma}^2)+2k+\frac{2k(k+1)}{n-k-1},
\qquad \widehat{\sigma}^2=\frac{1}{n}\sum_i (y_i-\hat y_i)^2,
\end{align}
with $k$ the number of free parameters (nonzero coefficients $+$ intercept). We also track $R^2$ and adjusted $R^2$:
\begin{align}
R^2=1-\frac{\sum_i(y_i-\hat y_i)^2}{\sum_i(y_i-\bar y)^2},\qquad
\widetilde{R}^2=1-(1-R^2)\frac{n-1}{n-k-1}.
\end{align}

\paragraph{Search.} A genetic algorithm evolves chromosomes using tournament selection, single‑point crossover, and per‑gene mutation, with elitism. Parsimony emerges from the LASSO penalty and the AICc term.
\end{document}
